


\subsection{Risk Capital Attribution and Risk-Adjusted Performance Measurement}

\subsubsection{Risk capital, economic capital, \& regulatory capital}

\begin{itemize}
\item Risk capital is the cushion that provides protection against the various risks (i.e., unexpected loss) inherent in the business of a corporation so that the firm can maintain its financial integrity and remain a going concern even in the event of a near-catastrophic worst-case scenario.
\item Risk capital is often called economic capital, and in most instances the generally accepted convention is that risk capital and economic capital are identical.
\item Regulatory capital
    \begin{itemize}
    \item First, regulatory capital applies to regulated industries to protect the interests of small depositors or policy holders.
    \item Second, regulatory capital is calculated by regulatory rules and formulas and sets a minimum required level of capital adequacy.
    \item Third, even if regulatory and risk capital are similar numbers, they may not be similar for each constituent business line.
    \end{itemize}
\end{itemize}



\subsubsection{Uses of economic capital-traditional way}
\begin{itemize}
\item Economic capital is really important to financial institutions.
    \begin{itemize}
    \item Capital is primarily used in a financial institution not only
    to provide funding but also to absorb risk.
    \item A bank's target solvency is a vital part of the product the
    bank is selling.
    \item Although bank creditworthiness is critical, banks are also
    highly opaque institutions.
    \item Banks operate in highly competitive financial markets.
    \end{itemize}
\end{itemize}



\subsubsection{Uses of economic capital-new way}
\begin{itemize}
\item Economic capital's new uses including:
    \begin{itemize}
    \item Performance measurement and incentive compensation at the firm, business unit, and individual levels
    \item Active portfolio management for entry/exit decisions
    \item Pricing transactions: calculate risk-based pricing for individual transactions. Each transaction considers the expected loss and the cost of economic capital allocated.
    \end{itemize}
\end{itemize}


\subsubsection{RAROC: risk-adjusted return on capital}
\begin{itemize}
\item \begin{equation}
    \text{RAROC} = \frac{\text{after-tax expected risk-adjusted net income}}{\text{economic capital}}
    \end{equation}
    \begin{itemize}
    \item After-tax expected risk-adjusted net income = expected revenues - costs - expected losses - taxes+ return on risk capital +/- transfers
    \item Expected revenues are the revenues that the activity is expected to generate (assuming no losses).
    \item Costs are the direct expenses associated with running the activity.
    \item Expected losses are also called PD amount, they correspond to the loan loss reserve (or provision) that the bank must set aside as the cost of doing business.
    \item Taxes are the expected amount of taxes imputed to the activity using the effective tax rate of the company.
    \item Return on risk capital is the return on the risk capital allocated to the activity. It is generally assumed that this risk capital is invested in risk-free securities.
    \item Transfers correspond to transfer pricing mechanisms, primarily between the business unit and the treasury group, such as charging the business unit for any funding cost incurred by its activities and any hedging cost.
    \end{itemize}
\item Economic capital is sum of risk capital and strategic capital
    \begin{itemize}
    \item Economic capital = risk capital+ strategic capital
    \item Strategic capital = goodwill + burned-out capital
    \item Economic capital = risk capital + goodwill + burned-out capital
    \item Strategic risk capital refers to the risk of significant investments about whose success and profitability there is high uncertainty.
    \item Burned-out capital refers to the idea that capital is spent on the initial stages of starting up business, but the business may ultimately fail.
    \item Goodwill is the amount paid above the replacement value of the net assets (assets - liabilities) when acquiring a firm.
    \end{itemize}
\end{itemize}


\subsubsection{Example}
\begin{itemize}
\item Claire bank has a 1 billion loans with 20\% return rate, the direct cost of this loan is 10\%, the expected loss is 5\%, and the return on government security is 2\%. This firm has 500 million economic capital. This firm has 25\% effective tax rate, compute the RAROC for the firm.
\end{itemize}


\subsubsection{Answer}
\begin{itemize}
\item Expected revenue = 20\% X 1,000 = 200 million
\item Direct cost= 10\% X 1,000 = 100 million
\item Expected losses = 5\% X 1,000 = 50 million
\item Return on risk capital = 2\% X 500 = 10 million
\item Transfers = 0
\item Economic capital = 500 million
\item \begin{equation}
    \text{RAROC} = \frac{(200 - 100 - 50 + 10) \times (1 - 25\%)}{500} = \frac{45}{500} = 9\%
    \end{equation}
\end{itemize}


\subsubsection{RAROC calculation issues}
\begin{itemize}
\item Time horizon choice for RAROC is somewhat subjective.
\item PD is used for economical capital calculation on an anticipatory or ex ante basis.
    \begin{itemize}
    \item A point-in-time (PIT) PD is reasonable for near-term expected losses and individual financial instrument.
    \item A through-the-cycle (TTC) PD is more reasonable for economic capital, current profitability, and strategic decisions.
    \item PIT is more volatile than TTC.
    \end{itemize}
\item Confidence level: setting a lower confidence level may reduce the amount of risk capital allocated to an activity.
    \begin{itemize}
    \item Lower confidence level → Lower economic capital →  Higher RAROC → Affect capital allocation decisions （allocating risk capital to business units and individual transactions for the purpose of measuring economic performance）
        \begin{itemize}
        \item Once capital has been allocated, each business unit manages its risk so that it does not exceed its allocated capital.
        \end{itemize}
    \end{itemize}
\item The confidence level in economic capital calculation should be consistent with firm's target credit rating and can be regarded as a quantitative expression of the risk appetite of the firm.
    \begin{itemize}
    \item Business units tend to apply lower confidence levels than the firm-wide target confidence level.
    \end{itemize}
\end{itemize}


\subsubsection{RAROC for capital budgeting decision}
\begin{itemize}
\item Hurdle rate ($h_{AT}$) is after-tax weighted average cost of equity capital to compare with RAROC to make investment decision.
\begin{equation}
    h_{AT} = \frac{CE \times R_{CE} + PE \times R_{PE}}{CE + PE}
    \end{equation}
$CE/R_{CE}$ and $PE/R_{PE}$ is the value/cost of common equity and preferred equity, respectively.
    \begin{itemize}
    \item If RAROC >hurdle rate, add value to the firm, the project should be accepted.
    \item If RAROC <hurdle rate, destroy value for the firm, the project should be rejected or closed down.
    \end{itemize}
\item ARAROC: Adjusted RAROC
    \begin{equation}
    ARAROC = RAROC - \beta_{CE} \left( \bar{R}_{M} - r_{f} \right)
    \end{equation}
\(\bar{R}_{M}\) is the expected return on the market portfolio, \(\beta_{CE}\) is the beta of the equity of the firm.
    \begin{itemize}
    \item ARAROC> $r_t$: Accept projects
    \item ARAROC< $r_t$: Reject projects
    \end{itemize}
\item Risk capital and diversification: risk capital for the firm should be significantly less than the sum of the stand-alone risk capital of the individual business units due to diversification.
    \begin{itemize}
    \item Stand-alone capital is the capital used by mutualindependent view among business units.
    \item Fully diversified capital is the capital attributed to each activity X and Y, considering diversification from combining.
    \item Marginal capital is the additional capital required by an incremental business, considering diversification.
    \end{itemize}
\end{itemize}


\subsubsection{Example}
\begin{itemize}
\item Assume a business unit in a bank plans to carry out only two
activities, X and Y:
\item Activity X alone requires \$40 of risk capital.
\item Activity Y alone requires \$50 of risk capital.
\item Activities X and Y combined need a total of \$70 of risk
capital.
\end{itemize}


\subsubsection{Answer}

\begin{table}[h]
    \centering
    \caption{Total Diversification Benefit}
    \renewcommand{\arraystretch}{1.5} % 设置行距为1.5

    \begin{tabular}{|c|c|c|}
        \hline
        \textbf{Total diversification benefit} & \textbf{X} & \textbf{Y} \\
        \hline
        Stand-alone capital & / & 40 \\
        \hline
        Fully-diversified capital & 20 & \( 40 - 20 \times \frac{40}{90} = 31.11 \) \\
        \hline
        Marginal capital & / & \( 70 - 50 = 20 \) \\
        \hline
        & \( 70 - 40 = 30 \) & \\
        \hline
    \end{tabular}
\end{table}








\subsection{Range of Practices and lssues in Economic Capital Frameworks}

\subsubsection{Overview of economical capital framework}
\begin{itemize}
\item Economic capital has broadened to combine the underlying risks (building blocks) into an overall aggregation. The building blocks:
    \begin{itemize}
    \item Use of economic capital and governance
    \item Risk measures
    \item Risk aggregation
    \item Validation of economic capital
    \item Dependency modelling in credit risk
    \item Counterparty credit risk
    \item Interest rate risk in the banking book
    \end{itemize}
\end{itemize}


\subsubsection{Issues in building blocks: governance}
\begin{itemize}
\item Senior management need to provide credible commitment, recognize the importance of using economic capital measures, put adequate resources and fully understand the limitations of economic capital measures.
\end{itemize}


\subsubsection{Issues in building blocks: risk measure}
\begin{itemize}
\item Banks use a variety of risk measures such as VaR and ES.
\item All risk measures have advantages and disadvantages which need to be understood within the context of their intended application.
\end{itemize}


\subsubsection{Issues in building blocks: risk aggregation}
\begin{itemize}
\item Practices and techniques in risk aggregation are generally less sophisticated than the methodologies that are used in measuring individual risk components.
\item Since individual risk components are typically estimated without much regard to the interactions between risks (e.g., between market and credit risk), the aggregation methodologies used may underestimate overall risk even if "no diversification" assumptions are used.
\item There are five commonly used aggregation methodologies.
    \begin{itemize}
    \item Simple summation: assumes equal weighting, does not
    take into account the diversification benefits （p = 1）.
    \item Constant diversification: subtracts a fixed diversification percentage from the overall amount.
    \item Variance-covariance matrix: Estimates of inter-risk correlations are difficult and costly to obtain. It does not adequately capture nonlinearities and skewed distributions, leading to underestimate the risk.
    \item Copulas: combines marginal probability distributions into a joint probability distribution.
        \begin{itemize}
        \item It offer greater flexibility in the aggregation of risks and promise a better approximation of the true risk distribution.
        \item More demanding input and building a joint distribution is very difficult.
        \end{itemize}
    \item Full modeling/simulation: simulates the impact of common risk drivers on all risk types and construct the joint distribution of losses.
        \begin{itemize}
        \item Theoretically, the most appealing method, potentially the most accurate method.
        \item High information technology demands, the process is time consuming, and it may provide a false sense of security.
        \end{itemize}
    \end{itemize}
\end{itemize}


\subsubsection{Issues in building blocks: validation}
\begin{itemize}
\item The validation of economic capital models is at a very preliminary stage. There exists a wide range of validation techniques, each of which provides evidence for （or against） only some of the desirable properties of a model.
\end{itemize}


\subsubsection{Issues in building blocks: dependency modeling}
\begin{itemize}
\item A particularly important and difficult aspect of portfolio credit the dependency structure, both linear and non-linear relationships between obligors.
\item The main concern in the area continues to center on the accuracy and stability of correlation during times of stress.
\item The correlation estimates provided by current models still depend heavily on explicit or implicit model assumptions.
\end{itemize}


\subsubsection{Issues in building blocks: counterparty credit risk}
\begin{itemize}
\item Measurement of counterparty credit risk represents a complex exercise, as it involves gathering data from multiple systems, millions of transactions, time horizons ranging from overnight to thirty or more years, tracking collateral and netting arrangements, thousands of counterparties.
\item This complexity creates unique market-risk-related challenges and operational risk-related challenges.
\end{itemize}


\subsubsection{Issues in building blocks: interest rate risk}
\begin{itemize}
\item The main challenges for interest rate risk in the banking book is the long holding period for balance sheet and modeling indeterminate cash flows due to embedded optionality in many banking book items.
\item If not adequately measured and managed, the asymmetrical payoff characteristics of instruments with embedded option features can present risks that are significantly greater than the risk measures suggest.
\end{itemize}


\subsubsection{Recommendations}
\begin{itemize}
\item Based on BIS recommendations, there are 10 aspects:
    \begin{itemize}
    \item 1.Use economic capital model in assessing capital adequacy
        \begin{itemize}
        \item A bank demonstrates how economic capital model is integrated into business decision-making process and understanding of gap between gross (standalone) and net ( diversified) risk.
        \end{itemize}
        
    \item 2.Senior mgt: credit commitment to ensure meaningfulness /integrity of process and adequate resources.
    \item 3.Transparency and integration into decision making
    \item  4.Risk identification: not all risks can be directly quantified.
    \item 5.Risk measures: no singularly preferred risk measure.
    \item 6.Risk aggregation: depends on quality of the measure of individual risk components and interactions between risks.
    \item 7. Validation: Validation should ensure the capital generated by the model is sufficient and model is fit for purpose.
    \item 8. Dependency modeling in credit risk: using adequate supplementary approaches （sensitivity analysis, scenario analysis） to address the limitation of credit models.
    \item 9.Counterparty credit risk: complementary measurement processes such as stress testing should be used.
    \item 10.lnterest rate risk in banking book: close attention to the instruments with embedded option, which can present significant risk.
    
    
    \end{itemize}
\end{itemize}


\subsubsection{Benefits and impacts of EC using and governance}
\begin{itemize}
\item In order to achieve a common measure across all risks and businesses, economic capital is often parameterized as an amount of capital that a bank needs to absorb unexpected losses over a certain time horizon at a given confidence level.
\item Because expected losses are accounted for in the pricing of a bank's products and loan loss provisioning, it is only unexpected losses that require economic capital.
\end{itemize}


\subsubsection{Benefits and impacts of EC using}
\begin{itemize}
\item Credit portfolio management
    \begin{itemize}
    \item Benefit
        \begin{itemize}
        \item Economic capital is a measurement of the level of concentration. A credit's incremental risk contribution is used to determine the concentration of the credit portfolio.
        \item The bank allocates economic capital according to the credit asset's marginal contribution to the bank's total portfolio risk rather than its standalone risk.
        \end{itemize}
    \item Impact
        \begin{itemize}
        \item Credit portfolio management allows banks to specify appropriate hedging strategies to reduce portfolio concentration, protecting against risk deterioration
        \end{itemize}
    \end{itemize}
\item Risk-based pricing
    \begin{itemize}
    \item Benefit
        \begin{itemize}
        \item In theory, under the assumption of competitive markets, prices are exogenous to banks acting as price-takers.
        \item Decisions are driven by a floor/threshold （the minimum RAROC） computed according to the amount of economic capital allocated to the deal.
        \end{itemize}
    \item Impact
        \begin{itemize}
        \item Exceptions (specific customer relationship/reputational side-effect) to the rule are strictly monitored and require the decision be elevated to a higher level of management.
        \end{itemize}
    \end{itemize}
\item Customer and product profitability analysis
    \begin{itemize}
    \item Benefit
        \begin{itemize}
        \item The measurement of performance can be extended down to the customer level, through the analysis of customer profitability to provide the relative riskadjusted profitability of customer relationships and
        products.
        \end{itemize}
    \item Impact
        \begin{itemize}
        \item Segmenting customers and resources re-allocating based on more efficient profitable relationships.
        \end{itemize}
    \end{itemize}
\item Management incentives
    \begin{itemize}
    \item Benefit
        \begin{itemize}
        \item Economic capital could affect the objective functions of decision-makers at the business unit level and be used in the incentive structure for business-unit management.
        \end{itemize}
    \item Impact
        \begin{itemize}
        \item Incentives are the most sensitive element for a majority of bank managers and motive them involved in the technical aspects of the economic capital allocation process.
        \item RAROC can motivate business unit managers by enhancing their incentive compensation as less economic capital is allocated to their respective business unit.
        \end{itemize}
    \end{itemize}
    
\end{itemize}


\subsection{Capital Planning at large Bank Holding Companies}



\subsubsection{Overview of Fed capital plan rule}
\begin{itemize}
\item Federal Reserve's Capital Plan Rule requires all U.S.BHCs to have a capital plan by a robust capital adequacy process. The capital plan covers all U.S. domiciled BHCs with total consolidated assets equal to \$50 billion or more.
\item Principle 1: Sound foundational risk management
\item Principle 2: Effective loss-estimation methodologies
\item Principle 3: Solid resource-estimation methodologies
\item Principle 4: Sufficient capital adequacy impact assessment
\item Principle 5: Comprehensive capital policy and capital planning
\item Principle 6: Robust internal controls
\item Principle 7: Effective governance
\end{itemize}


\subsubsection{Foundational risk management}
\begin{itemize}
\item The BHC has a sound risk-measurement and riskmanagement infrastructure to support identification, measurement, assessment and control of all material risks.
    \begin{itemize}
    \item Stronger risk-identification practices include standardized processes (regularly update, review risk exposures).
    \item Lagging practice: did not try to account for reputational risk, strategic risk and compliance risk.
    \end{itemize}
\end{itemize}


\subsubsection{Loss-estimation methodologies}
\begin{itemize}
\item The BHC has processes for translating risk measures into estimates of losses over a range of stressful scenarios and for aggregating those estimated losses across the BHC.
    \begin{itemize}
    \item Apply conservative approaches to estimate losses.
    \item Stronger scenario-design practices have clear narrative describing the vulnerabilities/material risks the firm facing.
    \item Stronger scenario-design practices used internal models in combination with expert judgment rather than relying solely on either.
    \item BHCs can use third-party scenarios with stronger practices tailoring third party-defined scenarios to their own risk profiles and unique vulnerabilities.
    \item Leading practice: In evaluating modeling program of BHC, BHC should strive to estimate loss models at granular level.
    \item BHCs should use internal operational-loss data as a starting point to provide historical perspective, and then incorporate forward-looking elements, idiosyncratic risks, and tail events to estimate losses.
    \item Lagging practice: model the loss given default of the bank's loan portfolio using a weighted-average approach at the aggregate portfolio level. Leading practice clearly tied LGD to underlying risk drivers.
    \item BHCs with leading practices used automated processes, while BHCs with lagging practices exhibited a high degree of manual intervention in the aggregation process.
    \item By design, estimates based on long-run average behavior, including periods of economic expansion and downturn, are not appropriate for projecting losses under stress.
    \item It is weak to combine two models: roll-rate model and net charge-off model.
    \end{itemize} 
\end{itemize}


\subsubsection{Resource-estimation methodologies}
\begin{itemize}
\item The BHC has a clear definition of available capital resources.
    \begin{itemize}
    \item Quantitative basis and qualitative projections (expert judgment);
    \item Conservatism and credibility;
    \item Documentation.
    \end{itemize}
    
\end{itemize}


\subsubsection{Capital policy and capital planning}
\begin{itemize}
\item BHC has capital policy for establishing capital goals, capital levels and composition of capital, making decisions about capital actions, and setting capital contingency plans.
    \begin{itemize}
    \item BHC should establish capital targets above capital goals.
    \item Contingency actions should be flexible enough and realistic for what is achievable during periods of stress.
    \item Lagging practice: establish triggers based on actual results but not on projected results, or based on minimum regulatory capital ratios.
    \end{itemize}
\end{itemize}


\subsubsection{Assessing capital adequacy impact}
\begin{itemize}
\item The BHC has processes for estimating of losses and capital resources to assess impacts on capital adequacy.
    \begin{itemize}
    \item BHCs should have a process for projections of the size and composition of on/off-balance sheet and RWA.
    \item Balance projections are a key input to enterprise-wide scenario analysis given direct impact on the estimation of losses, PPNR (pre-provision net revenue), and RWA.
    \item Use a full-revaluation approach to model the bank's nonlinear positions (such as complex derivative exposures).
    \end{itemize}  
\end{itemize}



\subsubsection{Internal controls}
\begin{itemize}
\item The BHC has robust internal controls governing capital adequacy process components, including policies and procedures; change control; model validation and independent review; comprehensive documentation; and review by internal audit.
    \begin{itemize}
    \item BHCs should ensure that validation covers all models and assumptions used for capital planning purposes, including any adjustments management has made to the model estimates (management overlay).
    \end{itemize}
\end{itemize}



\subsubsection{Governance}
\begin{itemize}
\item The BHC has effective board and senior management oversight of the capital adequacy plan ( CAP).
    \begin{itemize}
    \item Board has final oversight accountability for capital planning and should receive related information at least quarterly.
    \item Senior management is responsible for ensuring that capital planning activities authorized by the board are implemented in a satisfactory manner and is accountable to the board for the effectiveness of those activities.
    \end{itemize}
\end{itemize}



\subsection{Stress Testing Banks}


\subsubsection{Stress testing - historical evolution}
\begin{itemize}
\item Supervisory capital assessment program （SCAP） is a US bank stress test in 2009.
\item A successful macro-prudential stress testing program, also the first macro-prudential stress test after the 2007-2008 financial crisis.
    \begin{itemize}
    \item Macro-prudential focuses on the soundness of the banking system as a whole （systematic risks）, but micro-prudential focuses on the soundness of the individual institution.
    \end{itemize}
\item The 2009 SCAP tested macro scenarios with three main variables: growth in GDP, unemployment rate, and the house price index (HPI).
\item SCAP has at least two components:
    \begin{itemize}
    \item Assessment of institutions: capital strength vs. capital "hole" that needs to be filled
    \item A credible way of filling that hole
    \end{itemize}
\item Features of stress testing, pre- and post-SCAP
\begin{table}[h]
    \centering
    \caption{Comparison of Pre-SCAP and Post-SCAP}
    \renewcommand{\arraystretch}{1.5} % 设置行距为1.5
    \begin{tabular}{|c|c|c|}
        \hline
        \textbf{Pre-SCAP} & \textbf{Post-SCAP} \\
        \hline
        Mostly single shock & Broad macro scenario and market stress \\
        \hline
        Product or business unit level & Comprehensive, firm-wide \\
        \hline
        Static & Dynamic and path dependent \\
        \hline
        Not usually tied to capital adequacy & Explicit post-stress common equity threshold \\
        \hline
        Losses only & Losses, revenues and costs \\
        \hline
    \end{tabular}
\end{table}
\item Before 2011 all supervisory stress tests applied the same scenarios on all banks (i.e., one-size-fits-all approach).
\item Mechanically uniform approaches pose problems, so the 2011 and 2012 CCAR (Comprehensive capital analysis and review) asked banks to submit results from banks' own stress scenarios (baseline and stress) besides the supervisory stress scenario in order to bring banks' own vulnerabilities to light.
\item 2011 CCAR required that only the macro scenario was published (disclosed), but no bank level results.
\item 2012 CCAR disclosed nearly the same level of detail as the 2009 SCAP, namely bank-level loss rates and dollar losses by major regulatory asset class.
    \begin{itemize}
    \item Disclosure was an important trait of the 2009 SCAP.
        \begin{itemize}
        \item Before 2009 SCAP: disclosed realized losses
        \item After 2009 SCAP: disclosed projected losses
        \end{itemize}
    \end{itemize}
\item 2011 European banking authority （EBA） conducts European stress test.
\item The EBA test includes both a baseline scenario and a stress scenano.
    \begin{itemize}
    \item The EBA also required significant disclosures. The disclosure was more extensive, including information on exposure by asset class and by geography. All bank level results were available to download in spreadsheet form.
    \end{itemize}
\end{itemize}



\subsubsection{Stress testing-challenges in designing scenarios}
\begin{itemize}
\item Principal challenges in designing stress scenarios is coherence, especially inherently multi-risk factors, making it more difficult to design a coherent stress test.
    \begin{itemize}
    \item The supervisor's key challenge is to state the joint outcomes of inherently multi-risk factors.
    \item Not every thing becomes terrible simultaneously.
        \begin{itemize}
        \item Currency depreciation vs. currency appreciation, riskier asset vs. safe haven assets (flight to quality).
        \end{itemize}
        
    \end{itemize}
\end{itemize}



\subsubsection{Stress testing-challenge in modeling losses and revenues}
\begin{itemize}
\item From macro-scenario to micro-outcome: losses/revenues, the first task is to map from the few macro-factors into the many intermediate risk factors that drive losses.
\item In a stress model, the starting balance sheet generates first quarter's income and loss, which in turn determines quarterend balance sheet. The bank must maintain its capital and liquidity ratios level during all quarters of the stress test time horizon, typically for two years.
\end{itemize}



